\section{Core idea}
The project separates two sources of variation that a frozen-model workflow can
confound:
\begin{equation}
  \text{measurable operating variation} \neq
  \text{persistent client heterogeneity}.
\end{equation}
For vehicle dynamics, longitudinal speed $\rho=v_x$ is measurable and varies
within a drive. Structural parameters such as mass, yaw inertia, and tire
cornering stiffness remain fixed for a client in the initial study. Client $i$ is
represented as
\begin{equation}
  x_{i,k+1}=A_i(\rho_k)x_{i,k}+B_i(\rho_k)u_{i,k},
\end{equation}
and is assumed to belong approximately to one persistent dynamics family $c(i)$:
\begin{equation}
  A_i(\rho)\approx A_{c(i)}(\rho),\qquad
  B_i(\rho)\approx B_{c(i)}(\rho).
\end{equation}
The central hypothesis is that LPV scheduling should explain operating-point
variation, whereas clustered federated learning should handle unknown persistent
structural variation. In geometric terms, the goal is to learn families of model
manifolds rather than clusters of unrelated frozen-model points.

\subsection{Research question and hypothesis}
\paragraph{Research question.}
Can federated system identification disentangle measurable operating-condition
variation from persistent fleet heterogeneity, enabling a heterogeneous fleet to
be represented and controlled by a small number of reusable LPV model and
controller families?

\paragraph{Hypothesis.}
Explicit LPV scheduling combined with structural specialization yields better
generalization and lower model/controller proliferation than either one global
LPV model or purely clustered LTI models.

